% Top reaction force vs. loading-pin displacement u_y = 8*time (mm), with and without RaM.
% The x range stops at the end of the last RaM segment (_6.csv, t = 0.35625 -> u_y = 2.85 mm).
% RaM segments _1 to _6 start with an all-zero row at the restart time, so that row is dropped.
\pgfplotsset{
    fcskipfirst/.style={x filter/.expression={\coordindex==0 ? nan : x}},
}
\begin{figure}[htbp]
    \centering
    \begin{tikzpicture}
        \begin{axis}[
                width=0.6\textwidth,
                height=0.4\textwidth,
                xmin=0, xmax=2.85,
                restrict x to domain=0:2.85,
                xlabel={Loading pin displacement $u_y$ (\unit{\milli\meter})},
                ylabel={Top reaction force},
                legend pos=north east,
                legend cell align=left,
                grid=major,
            ]
            \addplot[blue, line width=1.6pt] table[col sep=comma, x expr=8*\thisrow{time}, y=reaction_top] {results/data/frac_challenge/frac_challenge_1e4meshnotch_single_0.csv};
            \addlegendentry{No RaM}
            \addplot[red, line width=0.8pt] table[col sep=comma, x expr=8*\thisrow{time}, y=reaction_top] {results/data/frac_challenge/frac_challenge_meshnotch_1e4_0.csv};
            \addlegendentry{RaM}
            \foreach \i in {1,...,6} {
                    \addplot[red, line width=0.8pt, fcskipfirst, forget plot] table[col sep=comma, x expr=8*\thisrow{time}, y=reaction_top] {results/data/frac_challenge/frac_challenge_meshnotch_1e4_\i.csv};
                }
        \end{axis}
    \end{tikzpicture}
    \caption{Top reaction force vs.\ loading pin displacement with and without RaM.}\label{fig:challenge_reaction}
    \label{fig:sfc_reactiondisp}
\end{figure}
